


For supervised learning, transformer-based models trained at scale have shown impressive abilities to perform tasks given an input context, often referred to as few-shot prompting or in-context learning~\citep{brown2020language}. In this setting, a pretrained model is presented with a small number of supervised input-output examples in its context, and is then asked to predict the most likely completion (i.e. output) of an unpaired input, without parameter updates. Over the last few years in-context learning has been applied to solve a range of tasks~\citep{min2022rethinking} and a growing number works are beginning to understand and analyze in-context learning for supervised learning~\citep{chan2022data,garg2022can,razeghi2022impact,akyurek2022learning}.
In this work, our focus is to study and understand in-context learning applied to sequential decision-making, specifically in the context of reinforcement learning (RL) settings. Decision-making (e.g. RL) is considerably more dynamic and complex than supervised learning. Understanding and leveraging in-context learning here could potentially unlock significant improvements in an agent's ability to adapt and make few-shot decisions in response to observations from the world. Such capabilities are instrumental for practical applications ranging from robotics to recommendation systems. 

For in-context decision-making~\citep{laskin2022context,xu2022prompting,xu2023hyper}, rather than input-output tuples, the context takes the form of state-action-reward tuples representing a dataset of interactions with an unknown environments. The agent must leverage these interactions to understand the dynamics of the world and what actions lead to good outcomes. A hallmark of good decision-making in online RL algorithms is a judicious balance of selecting exploratory actions to gather information and selecting increasingly optimal actions by exploiting that information~\citep{sutton2018reinforcement}. In contrast, an RL agent with access to only a suboptimal offline dataset should produce a policy that conservatively selects actions~\citep{levine2020offline}. An ideal in-context decision-maker should exhibit similar behaviors.

To study in-context decision-making formally, we propose a new simple supervised pretraining objective, namely, to train (via supervised learning) a transformer to predict an optimal action label\footnote{If not explicitly known, the optimal action can be determined by running any (potentially inefficient) minimax-optimal regret algorithm for each pretraining task.} given a query state and an in-context dataset of interactions, across a diverse set of tasks. 
We refer to the pretrained model as a Decision-Pretrained Transformer (DPT). Once trained, DPT can be deployed as either an online or offline RL algorithm in a new task by passing it an in-context dataset of interactions and querying it for predictions of the optimal action in different states. For example, online, the in-context dataset is initially empty and DPT's predictions are uncertain because the new task is unknown, but it fills the dataset with its interactions as it learns and becomes more confident about the optimal action. We show empirically and theoretically that DPT yields a surprisingly effective in-context decision-maker with regret guarantees. As it turns out, DPT effectively performs posterior sampling --- a provably sample-efficient Bayesian RL algorithm that has historically been limited by its computational burden~\citep{osband2013more}. We summarize our main findings below.
\begin{itemize}[leftmargin=1em]
\item \textbf{Predicting optimal actions alone gives rise to near-optimal decision-making algorithms.}
The DPT objective is solely based on predicting optimal actions from in-context interactions. At the outset, it is not immediately apparent that these predictions at test-time would yield good decision-making behavior when the task is unknown and behaviors such as
%some amount of 
online exploration are necessary to solve it.
Intriguingly, DPT as an algorithm is capable of dealing with this uncertainty in-context. For example, despite not being explicitly trained to explore, DPT exhibits an exploration strategy on par with hand-designed algorithms, as a means to discover the optimal actions. 


\item \textbf{DPT generalizes to new decision-making problems, offline and online.} We show DPT can handle reward distributions unseen in its pretraining data on bandit problems as well as unseen goals, dynamics, and datasets in simple MDPs. This suggests that the in-context strategies learned during pretraining are robust and generalizable without any parameter updates at test time.

\item \textbf{\macroName{} improves over the data used to pretrain it by exploiting latent structure.}
As an example, in parametric bandit problems, specialized algorithms can leverage structure (such as linear rewards) and offer provably better regret, but a representation must be known in advance. Perhaps surprisingly, we find that pretraining on linear bandit problems, even with unknown representations, leads DPT to select actions and explore in a way that matches an efficient linear bandit algorithm. This holds even when the source pretraining data comes from a suboptimal algorithm (i.e., one that does not take advantage of any latent structure), demonstrating the ability to learn improved in-context strategies beyond what it was trained on.  

\item \textbf{Posterior sampling can be implemented via in-context learning.} 
Posterior sampling (PS), a generalization of Thompson Sampling, can provably sample-efficiently solve online RL problems~\citep{osband2013more}, but a common criticism is the lack of computationally efficient ways to update and sample from a posterior distribution. \macroName{} can be viewed as learning a posterior distribution over optimal actions, shortcutting the PS procedure. Under some conditions, we show theoretically that \macroName{}  in-context is equivalent to PS. Furthermore, \macroName{}'s prior and posterior updates are grounded in data rather than needing to be specified \textit{a priori}. This suggests that in-context learning could help unlock practical and efficient RL via posterior sampling.

\end{itemize}

